% =========================================================================
% APPENDICE B — PROTOCOLLI DELLO STUDIO
% =========================================================================

\chapter{Protocolli dello studio}
\label{app:protocols}

Questa appendice documenta gli strumenti della valutazione empirica. Lo studio è erogato come questionario online autoguidato e auto-somministrato, perciò il protocollo è il questionario stesso: la sua struttura, le sue quattro varianti, i suoi item di comprensione, il testo di consenso e informazione che lo precede, e la griglia di codifica rispetto a cui le risposte sono analizzate. Il testo completo delle quattro varianti del questionario è riprodotto nella \S\ref{app:b-fulltext}; la chiave analitica che governa come le loro risposte aperte sono lette è esposta nella \S\ref{sec:app-coding-grid}.

% -------------------------------------------------------------------------
\section{Varianti del questionario}
\label{app:b-variants}

Lo strumento esiste in quattro varianti, incrociando due lingue con due condizioni:

\begin{itemize}
    \item \textbf{Lingua:} italiano e inglese, così che i partecipanti rispondano nella propria prima lingua. Le risposte aperte sono analizzate nella lingua in cui sono scritte (Appendice~\ref{app:limits-evaluate}).
    \item \textbf{Condizione:} una versione di \emph{controllo} che mostra solo la rappresentazione statica convenzionale (C1), e una versione \emph{sperimentale} che mostra la rappresentazione dinamica ed esperienziale (C2 e C3).
\end{itemize}

\noindent Le quattro varianti (IT-controllo, IT-sperimentale, EN-controllo, EN-sperimentale) sono identiche nella struttura, nelle domande di background e negli item di comprensione. Differiscono in esattamente quattro punti: la sezione \textsc{cosa farai}, i tre blocchi di esplorazione in apertura di scena, il paragrafo dei controlli della Scena~III e l'etichetta di rivelazione del gruppo. Quei quattro punti sono esposti affiancati nella Tabella~\ref{tab:b-differences}, e tutto il resto è riprodotto nella \S\ref{subsec:b-shared}. Ogni altra sezione è mantenuta costante, così che i due gruppi ricevano le stesse istruzioni e gli stessi item e differiscano solo nella forma di rappresentazione delle scene (Appendice~\ref{app:limits-evaluate}).

% -------------------------------------------------------------------------
\section{Struttura del questionario}
\label{app:b-structure}

Ogni partecipante segue lo stesso percorso, in quattro fasi.

\paragraph{1. Consenso e informazione.}
Il questionario si apre con la scheda informativa in linguaggio piano e la dichiarazione di consenso (\S\ref{app:b-consent}). La partecipazione procede solo dopo che il consenso è stato dato.

\paragraph{2. Domande di background.}
Un breve insieme di domande caratterizza il partecipante e permette di tenere conto della conoscenza pregressa nell'analisi:
\begin{itemize}
    \item fascia d'età;
    \item titolo di studio più elevato;
    \item ambito di studio o lavoro;
    \item familiarità auto-valutata con la scienza del clima, con l'AMOC in particolare, e con la lettura di visualizzazioni di dati.
\end{itemize}
Questi item non raccolgono alcuna informazione direttamente identificativa; le risposte sono anonimizzate (Appendice~\ref{app:ethics-conduct}).

\paragraph{3. Esposizione guidata con comprensione per scena.}
Il partecipante è condotto attraverso le tre scene della condizione assegnata in sequenza. Ogni scena si apre con un'introduzione neutra che porta la sua domanda guida, i link all'artefatto assegnato, e al ritorno pone le sue domande di comprensione prima che inizi la scena successiva, così che la comprensione sia catturata per ciascuna scena mentre si forma anziché ricostruita alla fine. Gli item di comprensione sono descritti nella \S\ref{app:b-comprehension}.

\paragraph{4. Riflessione finale e debriefing.}
Dopo le tre scene, una breve riflessione chiede al partecipante di fare un passo indietro e considerare il sistema nel suo insieme. Solo allora il debriefing rivela la struttura a due condizioni e svela in quale condizione il partecipante si trovava. A questo punto a ogni partecipante --- incluso il gruppo di controllo --- è dato accesso all'insieme completo dei materiali, l'altra condizione inclusa. Questa fase conclusiva non è valutata e non entra nel confronto tra gruppi; funge da svelamento e da restituzione al partecipante, ed è collocata per disegno fuori dal percorso valutato (Appendice~\ref{app:limits-evaluate}).

% -------------------------------------------------------------------------
\section{Item di comprensione}
\label{app:b-comprehension}

Gli item di comprensione sono l'evidenza primaria dello studio. Ogni scena pone due domande aperte di descrizione, un'auto-valutazione chiusa della fiducia e un prompt aperto finale per qualsiasi altra cosa il partecipante desideri aggiungere. Le domande aperte prendono di mira la proprietà strutturale che ciascuna scena è progettata per trasmettere, così da misurare la comprensione di ciò che l'artefatto effettivamente codifica anziché la conoscenza climatica generale.

Il formato aperto è deliberato: rivela il modello mentale di un partecipante con parole proprie anziché testarne la capacità di riconoscere una risposta fornita. L'item chiuso di fiducia è chiesto solo dopo la descrizione libera a cui si riferisce, mai prima, così che valutare la propria certezza non possa rimodellare il resoconto (\S\ref{sec:app-coding-grid}). Le risposte aperte sono il materiale primario per l'analisi; le valutazioni di fiducia sono lette solo in incrocio con il fatto che la struttura sia stata colta. Nessuna lettura strutturale è mai enunciata al partecipante in alcuna delle due condizioni: lo strumento insegna come leggere ciascuna scena senza fornire ciò che mostra.

% -------------------------------------------------------------------------
\section{Testo completo delle varianti del questionario}
\label{app:b-fulltext}

Questa sezione riproduce lo strumento come somministrato. Le quattro varianti condividono un
unico corpo di testo; differiscono in quattro punti, e quei quattro punti sono esposti affiancati
nella \S\ref{subsec:b-difference} così che l'equivalenza informativa rivendicata nella
\S\ref{app:b-variants} possa essere verificata anziché presa sulla fiducia.

Tre convenzioni valgono ovunque. I formati di risposta sono dati tra parentesi quadre come note
editoriali anziché riprodotti come controlli di modulo, e la cornice di interfaccia che Google
Forms aggiunge a un questionario esportato è omessa. I titoli di sezione sono riprodotti nel
maiuscoletto che il modulo usava. Il testo è qui riprodotto nella variante italiana, come
somministrato ai partecipanti che hanno scelto l'italiano; le varianti inglesi erano una
traduzione diretta dello stesso testo, con gli stessi item nello stesso ordine, e non sono
riprodotte separatamente.\footnote{Le due esportazioni inglesi differiscono in un piccolo numero
di dettagli tipografici accidentali --- qualche grafia corretta in una versione e non nell'altra.
Questi non toccano il contenuto o il significato di alcun item, e sono qui segnalati per
completezza anziché trattati come una differenza tra condizioni.}

% -------------------------------------------------------------------------
\subsection{Testo condiviso}
\label{subsec:b-shared}

\small

\paragraph{\textsc{Apertura}}
Ciao! Sono Enrico e questo studio fa parte della mia tesi magistrale in Sustainability, Innovation \&
Technology alla Tomorrow University (con un focus su Data Science e Machine Learning).

\emph{Living Atlantic} è il nome dell'insieme di visualizzazioni che vedrai.

Questo lavoro nasce da una domanda semplice da porre e difficile da risolvere: come si rende
comprensibile un sistema quando non lo si può vedere?

Molte parti del sistema climatico ci sono visibili --- il ghiaccio che si scioglie, i mari che si
alzano, le foreste che cambiano. Altre no. Alcune si muovono sotto la superficie, si dispiegano su
scale temporali più lunghe di una vita umana, e non possono essere comprese guardando un singolo
luogo o momento. Eppure possono essere altrettanto rilevanti --- specialmente quando un sistema si
avvicina a un punto di tipping, dove un cambiamento graduale può portare a spostamenti improvvisi e
potenzialmente irreversibili.

Il problema non è la mancanza di dati. Abbiamo dati in abbondanza. La sfida è trasformare quei dati
in comprensione. La mente umana può faticare a cogliere sistemi che fluiscono, accumulano,
interagiscono e cambiano nel tempo --- in particolare quando quelle dinamiche sono ridotte a un
grafico immobile o ad altre forme tradizionali e statiche di visualizzazione e comunicazione. Anche
persone molto capaci possono trovarlo difficile. Non è semplicemente una questione di quanto si sa.

È qui che questo studio esplora una possibilità diversa: \textbf{l'IA e il machine learning possono
sostenere una comprensione cognitiva più profonda dei sistemi fisici complessi, in particolare
quelli che coinvolgono punti di tipping? E possono aiutarci ad andare oltre il semplice vedere i
dati verso il comprendere davvero come un sistema si comporta?}

Questo studio esplora uno di questi approcci e ti chiede di aiutarmi a scoprire se funziona. Non ti
sarà chiesto di studiare nulla in anticipo, né di dare risposte ``giuste''. Mi interessa una cosa
sola: come leggi, interpreti e comprendi ciò che vedi.

Lo studio richiederà circa 20 minuti. Puoi interrompere in qualsiasi momento. La prima parte si
concentrerà su di te, e poi entreremo nel cuore del progetto. Grazie ancora per aver dedicato un
po' di tempo a questo!

\paragraph{\textsc{Consenso e anonimato}}
La partecipazione a questo studio è del tutto volontaria. Il questionario è anonimo, e non vengono
raccolti nomi, indirizzi email o altre informazioni che potrebbero identificarti direttamente.

Le tue risposte saranno conservate in modo sicuro e usate unicamente ai fini di questa ricerca.
Nessuna informazione personalmente identificabile sarà inclusa nell'analisi o riportata nei
risultati.

Sei libero di interrompere la partecipazione in qualsiasi momento, senza conseguenze e senza dover
fornire una ragione.

\textbf{Proseguendo con il questionario, confermi di aver compreso le informazioni sopra riportate e
di accettare volontariamente di partecipare allo studio.}

\paragraph{\textsc{Disclaimer su immagini / visualizzazioni}}
Tutte le immagini illustrative usate in questo questionario provengono da Unsplash o Google Images,
mentre tutte le visualizzazioni di dati sono state create appositamente per questo studio
dall'autore.

\paragraph{\textsc{Qualcosa su di te}}
Questa sezione pone alcune domande su di te e sul tuo background. Come prima, tutte le risposte
restano anonime. Queste domande sono incluse unicamente per sostenere l'analisi statistica dello
studio e per aiutarmi a comprendere e interpretare meglio i risultati.

\begin{enumerate}[leftmargin=2em, itemsep=3pt]
\item Qual è la tua fascia d'età? \emph{[scelta singola: 18--24; 25--34; 35--44; 45--54; 55--64;
      65+]}
\item Qual è il tuo titolo di studio più elevato? \emph{[scelta singola: Elementari; Medie;
      Scuola superiore; Laurea triennale; Laurea magistrale; Dottorato; Altro]}
\item Qual è il tuo ambito di studio / lavoro? (Sii conciso nella risposta).
      \emph{[testo libero]}
\item Quanto hai familiarità con la scienza del clima e con i temi legati al clima?
      (1 = per niente, 10 = esperto/a) \emph{[scala 1--10, ancorata ``Per niente'' /
      ``Esperto/a'']}
\item Quanto hai familiarità con la circolazione dell'Oceano Atlantico e in particolare con l'AMOC
      (Atlantic Meridional Overturning Circulation)? (1 = per niente, 10 = molto esperto/a)
      \emph{[scala 1--10, ancorata ``Per niente'' / ``Molto esperto/a'']}
\item Quanto hai familiarità nel leggere / interpretare grafici, mappe e visualizzazioni di dati
      (anche quando trattano argomenti che potresti non conoscere bene)? (1 = per niente,
      10 = molto competente) \emph{[scala 1--10, ancorata ``Per niente'' / ``Molto competente'']}
\end{enumerate}

\paragraph{\textsc{Come funziona e cosa stai per guardare}}
Dedica qualche minuto a leggere con attenzione le informazioni qui sotto. Forniscono il background
essenziale che ti servirà per comprendere ciò che vedi nelle scene che seguono --- incluso il
sistema esplorato, i concetti chiave e le unità usate per descriverlo.

Non devi studiare o memorizzare nulla. Tuttavia, leggi con attenzione e tieni a mente queste
informazioni mentre procedi attraverso le diverse scene, perché forniranno il contesto comune per
interpretare ciò che vedi.

\paragraph{\textsc{Caso di studio}}
Come caso di studio ho scelto l'AMOC (Atlantic Meridional Overturning Circulation). È un vasto
sistema di acqua marina in movimento che circola attraverso l'Oceano Atlantico, ridistribuendo
calore, sale e nutrienti nell'oceano e svolgendo un ruolo importante nel regolare il clima, in
particolare in Europa e nelle sue vicinanze.

L'AMOC non è qualcosa che possiamo vedere direttamente. Si estende su un intero oceano, si muove in
continuazione e opera su scale temporali difficili da esperire nell'arco di una vita umana. Dal
2004, una linea di strumenti a circa \ang{26.5}~nord ha monitorato con continuità l'Atlantico,
registrando quanta acqua si muove, a quale profondità, e come questi flussi cambiano nel tempo.

\paragraph{\textsc{Perché questo sistema?}}
Ho scelto l'AMOC come caso di studio per questa ricerca per una ragione semplice: riunisce molte
delle caratteristiche che rendono i sistemi fisici complessi difficili da comprendere per gli esseri
umani.

È invisibile, dinamico, distribuito nello spazio e in costante cambiamento. Il suo comportamento
emerge da molti processi interagenti, e può subire cambiamenti bruschi quando soglie critiche sono
avvicinate o superate. Non possiamo semplicemente guardare il sistema e capire cosa sta facendo.
Dobbiamo ricostruirne il comportamento da misurazioni, modelli e rappresentazioni visive.

Allo stesso tempo, l'AMOC non è solo un interessante esempio scientifico. È un sistema che conta.
Il suo comportamento è strettamente connesso al sistema climatico e a cambiamenti che possono
interessare regioni e società su larga scala. Eppure molte persone non ne hanno mai sentito
parlare, o sono consapevoli della circolazione oceanica solo in termini molto generali.

Questo rende l'AMOC particolarmente utile per questo studio: \textbf{è un sistema reale e di grande
portata, difficile da percepire direttamente, ma che può essere rappresentato attraverso dati e
visualizzazioni.} Se riusciamo a trovare modi per rendere un sistema come questo più comprensibile,
gli stessi principi potrebbero potenzialmente essere applicati a molti altri sistemi invisibili e
complessi che plasmano il nostro futuro.

\paragraph{\textsc{L'unità}}
L'intensità dell'AMOC si misura in Sverdrup (Sv). Uno Sverdrup è un milione di metri cubi d'acqua
che fluiscono al secondo. Quindi quando vedi un valore espresso in Sv, rappresenta il volume d'acqua
che attraversa il sistema ogni secondo.

\paragraph{\textsc{Cosa farai}}
\emph{[Questo blocco è uno dei quattro punti di differenza; vedi \S\ref{subsec:b-difference}.]}

\paragraph{\textsc{Nota rapida}}
Per questo studio, i partecipanti sono stati assegnati casualmente a uno di due modi diversi di
esplorare l'AMOC: una condizione a visualizzazione statica o una condizione
dinamica-interattiva-esperienziale.

Esperirai una di queste due condizioni. Quale ti sia stata assegnata sarà rivelato alla fine dello
studio.

Per ora, segui semplicemente le istruzioni ed esplora ciò che vedi nel modo più naturale possibile.
Non c'è nulla che tu debba sapere in anticipo --- e non c'è un modo giusto di esperirlo.

\paragraph{\textsc{Scena 1 --- La colonna}}
Stai per guardare la colonna d'acqua dell'Atlantico vista di taglio, dalla superficie fino al
fondale, migliaia di metri più in basso. La curva che vedrai descrive questa colonna.

\emph{[Domanda guida e invito all'azione: punto di differenza; vedi \S\ref{subsec:b-difference}. I
partecipanti aprivano poi la scena assegnata in una nuova scheda e tornavano al modulo.]}

\begin{enumerate}[leftmargin=2.2em, itemsep=3pt, start=7]
\item Descrivi con parole tue ciò che hai appena visto. Prova a descrivere come si muove l'acqua a
      qualcuno che non l'ha vista. Si muove tutta allo stesso modo? \emph{[testo libero]}
\item C'è stato qualcosa che ti ha sorpreso o incuriosito? Qualcosa di inatteso? Qualcosa che non
      hai colto affatto? \emph{[testo libero]}
\item Quanto sei sicuro/a delle tue descrizioni? (1 = per niente, 10 = molto sicuro/a)
      \emph{[scala 1--10, ancorata ``Per niente'' / ``Molto sicuro/a'']}
\item C'è altro che hai capito / che vorresti aggiungere? \emph{[testo libero]}
\end{enumerate}

\paragraph{\textsc{Scena 2 --- La sala macchine}}
In questa scena, ogni punto rappresenta lo stato dell'oceano in un dato mese: due punti vicini
indicano due mesi simili, mentre due punti distanti indicano due mesi diversi. Il centro rappresenta
l'equilibrio del sistema --- lo stato a cui l'oceano tende naturalmente a tornare quando se ne
allontana.

\emph{[Domanda guida e invito all'azione: punto di differenza.]}

\begin{enumerate}[leftmargin=2.4em, itemsep=3pt, start=11]
\item Descrivi con parole tue come si muove l'acqua nel tempo. Si muove a ritmo costante, o in modo
      irregolare? \emph{[testo libero]}
\item C'è stato un momento che sembrava diverso dagli altri? Qualcosa che ti ha colpito?
      \emph{[testo libero]}
\item Quanto sei sicuro/a delle tue descrizioni? (1 = per niente, 10 = molto sicuro/a)
      \emph{[scala 1--10]}
\item C'è altro che hai capito / che vorresti aggiungere? \emph{[testo libero]}
\end{enumerate}

\paragraph{\textsc{Scena 3 --- La soglia}}
In questa scena, guarderai tre diverse rappresentazioni dello stesso aspetto dell'AMOC, ciascuna
ricostruita usando un modello scientifico indipendente.

\emph{[Il paragrafo che descrive come le tre rappresentazioni possono essere esaminate, la domanda
guida e l'invito all'azione sono punti di differenza; vedi \S\ref{subsec:b-difference}.]}

\begin{enumerate}[leftmargin=2.4em, itemsep=3pt, start=15]
\item Descrivi con parole tue ciò che hai visto. Cosa stava succedendo tra i tre modelli? Che
      informazioni forniscono? \emph{[testo libero]}
\item I tre modelli concordano tra loro nel corso degli anni, oppure divergono in certi periodi?
      \emph{[testo libero]}
\item Quanto sei sicuro/a delle tue descrizioni? (1 = per niente, 10 = molto sicuro/a)
      \emph{[scala 1--10]}
\item C'è altro che hai capito / che vorresti aggiungere? \emph{[testo libero]}
\end{enumerate}

\paragraph{\textsc{Riflessione finale}}
Sei ora arrivato alla fine dell'esplorazione.

Prima di concludere, prenditi un momento per fare un passo indietro dalle singole scene e pensare
all'AMOC nel suo insieme. Qui non ci sono risposte giuste o sbagliate, e non sei messo alla prova su
ciò che hai imparato. Mi interessa che cosa, se qualcosa, senti di comprendere ora del sistema e come
sei arrivato a comprenderlo.

Rispondi alle seguenti domande con parole tue.

\begin{enumerate}[leftmargin=2.4em, itemsep=3pt, start=19]
\item \textsc{Comprensione} --- Cosa senti di aver ora compreso dell'AMOC che prima non capivi / non
      sapevi? \emph{[testo libero]}
\item \textsc{Chiarezza} --- C'è stato qualcosa, tra i concetti sull'AMOC presentati qui, che è
      rimasto per te confuso o poco chiaro? \emph{[testo libero]}
\item \textsc{Esperienza} --- Il modo in cui l'AMOC è stato presentato ha cambiato il tuo modo di
      pensare o comprendere il sistema? Se sì, come? \emph{[testo libero]}
\end{enumerate}

\paragraph{\textsc{Debriefing}}
\emph{[Mostrato solo dopo che tutte le sezioni precedenti erano state completate.]}

Grazie di cuore per aver preso parte. Il tuo modo di vedere, interpretare e descrivere il sistema è
una parte reale di questa ricerca. Ora che hai completato tutte le attività dello studio, posso
finalmente rivelare il quadro completo.

In totale 25 partecipanti hanno preso parte a questo studio. Dodici partecipanti sono stati
assegnati casualmente a un gruppo di controllo, che ha esperito l'AMOC esclusivamente attraverso
visualizzazioni tradizionali e statiche. I restanti partecipanti hanno esperito le rappresentazioni
dinamiche ed esperienziali esplorate in questo studio. Confrontare i due gruppi è una parte centrale
di questa ricerca: mi permette di indagare se il modo in cui un sistema è rappresentato possa
influenzare come le persone lo percepiscono, interpretano e comprendono.

Con un campione di queste dimensioni, lo studio non intende sostenere un'inferenza statistica robusta
o generalizzare i suoi risultati a una popolazione più ampia. Serve invece come indagine
esplorativa: cerca pattern emergenti, differenze nel modo in cui le persone interpretano il sistema,
e indicazioni su se questo approccio sia abbastanza promettente da giustificare ulteriori ricerche
con campioni più ampi.

Lungo tutto lo studio, gli stessi dati di fondo sull'AMOC sono stati presentati in forme diverse:
una rappresentazione statica tradizionale e una rappresentazione dinamica ed esperienziale. Le tre
scene si sono concentrate su aspetti diversi dello stesso sistema complesso:

\begin{itemize}[leftmargin=2em, itemsep=2pt]
\item \emph{La colonna} --- l'acqua scorre in direzioni opposte a profondità diverse: in generale
      verso nord vicino alla superficie e verso sud a maggiori profondità.
\item \emph{La sala macchine} --- il sistema ripassa per stati che ha già attraversato, rivelando
      aspetti del suo comportamento nel tempo.
\item \emph{La soglia} --- il livello di accordo tra modelli indipendenti cambia nel tempo, con
      maggiore disaccordo nei periodi più antichi e crescente convergenza man mano che le
      osservazioni migliorano.
\end{itemize}

Queste rappresentazioni sono state progettate per esplorare una domanda centrale della mia tesi:
cambiare il modo in cui esperiamo i dati può aiutarci a comprendere i sistemi fisici complessi in
modo più profondo?

Prosegui alla sezione successiva per scoprire a quale gruppo sei stato assegnato.

\paragraph{\textsc{Rivelazione del gruppo}}
\emph{[Punto di differenza; vedi \S\ref{subsec:b-difference}.]} Entrambe le varianti si chiudono con
la stessa affermazione: \emph{Le diverse condizioni non sono state pienamente svelate in anticipo
così che le aspettative dei partecipanti non influenzassero le loro risposte.}

\paragraph{\textsc{Esplora lo studio completo}}
Se desideri esplorare lo studio oltre la versione che hai esperito, puoi ora accedere a tutte e tre
le scene in tutte le rappresentazioni disponibili, incluse le versioni statica e
dinamica/esperienziale, cliccando sui seguenti link: \emph{[link alla Scena 1, Scena 2 e Scena 3,
ciascuna in tutte le rappresentazioni]}.

\paragraph{\textsc{Esplora tutti i materiali dello studio}}
Se sei interessato ad andare oltre lo studio, puoi esplorare il progetto, il suo codice e i
materiali di ricerca attraverso i link qui sotto. \emph{[Link ai due repository di ricerca ---
estrazione di pattern con machine learning e visualizzazioni --- documentati nell'Appendice~\ref{app:c-repos}.]}

La tesi completa, inclusi la metodologia di ricerca, il quadro teorico, il disegno sperimentale e
--- una volta disponibili --- i risultati di questo studio, sarà pubblicata qui: \emph{[link al PDF
della tesi]}. I risultati saranno resi disponibili non appena l'analisi sarà completa.

\paragraph{\textsc{Un ultimo grazie}}
Grazie ancora per aver dedicato il tempo a far parte di questa ricerca e per avermi aiutato a
indagare se un modo diverso di vedere un sistema complesso possa condurre a un modo diverso di
comprenderlo.

Se hai domande, riflessioni, o semplicemente vorresti saperne di più sul progetto, scrivimi pure a
\texttt{e.vaccari99@gmail.com}. Se sei curioso dei risultati, puoi consultare il file principale
della tesi tra circa un mese (set.\ 2026), dove ho in programma di pubblicare gli esiti di questo
studio.

\normalsize

% -------------------------------------------------------------------------
\subsection{Punti di differenza}
\label{subsec:b-difference}

Tutto ciò che è riprodotto sopra era identico per entrambi i gruppi. I quattro blocchi qui sotto sono
gli unici punti in cui il testo somministrato divergeva, e la divergenza è confinata al verbo di
ingaggio --- \emph{guardare} contro \emph{esplorare} --- e, nella Scena~III, alla descrizione di ciò
che l'artefatto assegnato permette di fare. Nessuna affermazione strutturale sull'AMOC, e nessuna
parte di alcun item di comprensione, differisce tra le due colonne.

\begin{table}[tbp]
\centering
\caption[Punti di differenza tra le varianti di controllo ed esperienziale]{I quattro punti in cui i
questionari di controllo ed esperienziale divergono. Tutto il resto --- l'apertura, la dichiarazione
di consenso e anonimato, gli item di background, il briefing su caso di studio e unità, gli item di
comprensione, la riflessione finale e il debriefing --- è comune a entrambi. La divergenza segue ciò
che ciascuna condizione permette di fare anziché ciò che afferma sul sistema.}
\label{tab:b-differences}
\footnotesize
\renewcommand{\arraystretch}{1.25}
\begin{tabularx}{\textwidth}{@{}p{2.3cm} X X@{}}
\toprule
\textbf{Blocco} & \textbf{Controllo (C1)} & \textbf{Esperienziale (C2+C3)} \\
\midrule

\makecell[l]{1.\\\textsc{cosa}\\\textsc{farai}} &
\emph{Prima, ti introdurrò brevemente la visualizzazione che stai per guardare}, senza darti troppi
dettagli. \dots{} \emph{Poi, aprirai la scena in una nuova scheda e la guarderai con attenzione}, per
tutto il tempo che vuoi. \dots{} Non c'è nulla da studiare in anticipo e nessuna risposta giusta.
\emph{Guarda attentamente, osserva,} e dimmi cosa noti. &
\emph{Prima, ti introdurrò rapidamente la visualizzazione che stai per guardare}, senza darti troppi
dettagli. \dots{} \emph{Poi, aprirai la scena in una nuova scheda e la esplorerai con calma}, per
tutto il tempo che vuoi. \dots{} Non c'è nulla da studiare in anticipo e nessuna risposta giusta.
\emph{Esplora, osserva,} e dimmi cosa noti. \\
\addlinespace

\makecell[l]{2.\\Invito all'azione\\in apertura di\\scena (tutte e\\tre le scene)} &
Tieni a mente questa domanda \emph{mentre guardi}: [domanda guida] \newline
\textsc{guarda attentamente!} \newline
Apri la scena in una nuova scheda cliccando il link qui sotto, \emph{esamina la figura con
attenzione}, poi torna e rispondi alle domande sotto. &
Tieni a mente questa domanda \emph{mentre esplori}: [domanda guida] \newline
\textsc{esplora!} \newline
Apri la scena in una nuova scheda cliccando sul seguente link, \emph{esplora con calma}, poi torna e
rispondi per favore alle domande sotto. \\
\addlinespace

\makecell[l]{3.\\Affordance\\della Scena~III} &
Le tre rappresentazioni \emph{sono mostrate} nel tempo, permettendoti di \emph{vedere} come il
sistema si comporta in diversi periodi. Prenditi il tempo per \emph{esaminare la figura} e presta
attenzione a come le tre rappresentazioni si relazionano tra loro. \emph{La figura mostra i tre
modelli insieme lungo l'intero periodo, così puoi confrontarli direttamente e vedere dove concordano
e dove differiscono.} &
Le tre rappresentazioni \emph{evolvono} nel tempo, permettendoti di \emph{osservare} come il sistema
si comporta in diversi periodi. Prenditi il tempo per \emph{esplorare la scena} e presta attenzione a
come le tre rappresentazioni si relazionano tra loro. \emph{Puoi passare da un modello all'altro,
mostrarli singolarmente o insieme, e regolare la finestra di lisciamento per guardare i dati su scale
temporali diverse: 3, 12 o 24 mesi. Questi controlli sono lì per lasciarti esplorare i dati da
prospettive diverse.} \\
\addlinespace

\makecell[l]{4.\\Rivelazione\\del gruppo} &
\textsc{facevi parte del gruppo di controllo!} Sei stato assegnato casualmente al gruppo di
controllo. Questo significa che hai esplorato l'AMOC attraverso visualizzazioni tradizionali e
statiche. Un altro gruppo di partecipanti ha esperito gli stessi dati di fondo attraverso una
rappresentazione interattiva e dinamica\dots &
\textsc{facevi parte del gruppo esperienziale!} Sei stato assegnato casualmente al gruppo
esperienziale. Questo significa che hai esplorato l'AMOC attraverso una rappresentazione interattiva
e dinamica\dots Potresti non averlo saputo, ma un altro gruppo di partecipanti ha esperito gli stessi
dati di fondo attraverso visualizzazioni tradizionali e statiche\dots \\

\bottomrule
\end{tabularx}
\end{table}

Due osservazioni seguono dalla tabella, ed entrambe riguardano l'affermazione di equivalenza della
\S\ref{app:b-variants}. Primo, i punti~1 e~2 sono lessicali: al gruppo di controllo si dice di
\emph{guardare} dove al gruppo esperienziale si dice di \emph{esplorare}, perché in una condizione
non c'è nulla su cui agire e nell'altra sì. Nessuna delle due formulazioni dice a un partecipante
alcunché sull'AMOC. Secondo, il punto~3 è l'unico luogo in cui i due testi descrivono affordance
diverse, e lo fa perché le affordance differiscono davvero: la Scena~III esperienziale porta
controlli che la figura statica non ha, e un partecipante che non fosse stato informato della loro
esistenza non avrebbe potuto usarli. Ciò che entrambe le versioni trattengono è identico --- nessuna
delle due nomina la lettura strutturale che la scena è costruita per trasmettere, che è ciò che rende
significativo il confronto nel Capitolo~\ref{ch:results}.

% -------------------------------------------------------------------------
\section{Scheda informativa e consenso}
\label{app:b-consent}

La scheda informativa, riprodotta come somministrata nella \S\ref{subsec:b-shared}, enunciava in linguaggio piano cosa comporta lo studio e quanto tempo richiede; che la partecipazione è volontaria e può essere interrotta in qualsiasi momento senza conseguenze; che non viene raccolta alcuna informazione direttamente identificativa; e che le risposte sono conservate in modo sicuro e usate solo per questa ricerca. L'indirizzo di contatto del ricercatore era dato in chiusura del questionario. Oltre a quanto la scheda enunciava, il ricercatore si è impegnato a trattare i dati in linea con il Regolamento generale sulla protezione dei dati, a conservarli solo per il tempo che la ricerca richiede, e a limitare l'accesso al ricercatore e al relatore. Il consenso è stato dato esplicitamente prima dell'inizio del questionario.


Gli impegni etici che questi documenti attuano sono esposti nell'Appendice~\ref{app:ethics-conduct}.

% =========================================================================
\section{La griglia di codifica}
\label{sec:app-coding-grid}

Le risposte aperte sono state analizzate rispetto alla griglia esposta sotto, fissata prima di leggere le risposte. Le sue categorie sono di due tipi. I \emph{target strutturali} registrano se un partecipante ha colto l'affermazione che ciascuna scena era costruita per consegnare, e sono riassunti nella Tabella~\ref{tab:coding-targets}. I \emph{codici diagnostici} registrano specifiche letture errate e proprietà trasversali della risposta, e sono definiti nel codebook che segue. Nulla nella griglia è mostrato ai partecipanti; governa l'analisi, non lo strumento.

La lettura è deduttiva al livello di queste categorie e interpretativa al loro interno: la griglia fissa ciò che si cerca, mentre il giudizio decide se una data descrizione, con parole proprie, soddisfi la definizione di una categoria. Ogni codice è applicato alle risposte aperte di ciascuna scena e mai all'item chiuso.

% -------------------------------------------------------------------------
\subsection{Target strutturali}
\label{subsec:app-targets}

Ogni scena porta un'affermazione strutturale, nominata dalla sua domanda guida e definita nella \S\ref{sec:three-scenes}. Una descrizione è codificata \emph{grasped} (colto) solo quando nomina quella struttura con parole proprie del partecipante; la presenza di movimento, colore o interesse generale non è sufficiente. Un credito intermedio è registrato come \emph{partial} (parziale) dove una componente necessaria è presente e un'altra assente, come specificato per ciascuna scena.

\begin{table}[htbp]
\centering
\caption[Target di codifica strutturale per scena]{Il target strutturale di ciascuna scena, con le condizioni sotto cui una descrizione libera è codificata \emph{grasped}, \emph{partial} o \emph{not grasped}. I target corrispondono uno a uno alle domande guida della \S\ref{sec:three-scenes}.}
\label{tab:coding-targets}
\small
\renewcommand{\arraystretch}{1.35}
\begin{tabularx}{\textwidth}{@{}p{2.5cm} X X@{}}
\toprule
\textbf{Scena} &
\textbf{Affermazione strutturale (target)} &
\textbf{Codificata \emph{grasped} quando\ldots} \\
\midrule

\makecell[l]{Scena I\\La colonna} &
L'acqua si muove in direzioni \emph{opposte} a profondità diverse: in generale verso nord vicino alla superficie, verso sud in profondità. &
la descrizione nomina (b) due direzioni opposte \emph{e} (c) lega la direzione alla profondità. \emph{Partial}: direzione menzionata (a) ma non sia opposta sia legata alla profondità. \\

\makecell[l]{Scena II\\La sala macchine} &
Il sistema \emph{ripassa} per stati che ha già occupato; il moto è ricorrente, non una deriva monotòna verso il nuovo. &
la descrizione trasmette ritorno, ciclicità o ripassaggio per stati precedenti. \emph{Partial}: moto o cambiamento notato senza l'idea di ritorno. \\

\makecell[l]{Scena III\\La soglia} &
L'\emph{accordo} tra i tre modelli \emph{cambia} nel tempo --- più disaccordo prima, convergenza dopo. &
la descrizione nomina un cambiamento nell'accordo nel tempo. \emph{Partial}: accordo o disaccordo notato come costante, senza cambiamento nel tempo. \\

\bottomrule
\end{tabularx}
\end{table}

% -------------------------------------------------------------------------
\subsection{Assi trasversali}
\label{subsec:app-axes}

Quattro assi sono registrati per ogni risposta, indipendentemente dal target strutturale, e sono la base per i conteggi descrittivi riportati nel Capitolo~\ref{ch:results}.

\begin{description}
\item[\textnormal{\textsc{Structure-grasped}}] Il verdetto dalla Tabella~\ref{tab:coding-targets}: grasped, partial o not grasped. L'esito primario, conteggiato per condizione.

\item[\textnormal{\textsc{Confidence}}] L'auto-valutazione del partecipante, registrata dall'item chiuso che segue ciascuna descrizione libera. Letta solo in incrocio, mai come esito in sé.

\item[\textnormal{\textsc{Confidence\,$\times$\,correctness}}] L'incrocio dei due precedenti. La cella d'interesse è alta fiducia con struttura \emph{non} colta --- la lettura sicura-ma-errata che \textcite{fischer2020} documentano, e che l'aspettativa~E3.2 riguarda.

\item[\textnormal{\textsc{Uncertainty-as-quality}}] Se il partecipante esprime incertezza \emph{produttivamente} --- nominando ciò di cui non è sicuro, o ciò che un display non gli dice --- come distinto sia dalla falsa fiducia sia dalla non-risposta vuota. L'incertezza produttiva è accreditata, non penalizzata; è la faccia positiva dello stesso asse la cui faccia negativa è la sicurezza-ma-errore.
\end{description}

% -------------------------------------------------------------------------
\subsection{Codici diagnostici}
\label{subsec:app-diagnostic}

Oltre al verdetto strutturale, la griglia registra una lettura errata, identificata in anticipo come probabile e dotata di peso interpretativo. È definita da ciò che autorizza il codice e da ciò che lo esclude, così che la sua applicazione sia ripetibile.

\paragraph{S1-MISC-INT --- Funzione di corrente letta come profilo di velocità.}
\emph{Si applica a: Scena~I.} Il partecipante legge l'\emph{altezza} della curva $\Psi(z)$ come la velocità o l'intensità del movimento dell'acqua a quella profondità, anziché leggere la sua \emph{pendenza} come il trasporto locale. Equivalentemente, $\Psi$ è trattata come un profilo di velocità diretto anziché come l'integrale cumulato del trasporto.

\emph{Indicatori positivi} (ne basta uno): velocità descritta come decrescente o crescente al passo con l'altezza della curva; ``nessun movimento'' o immobilità collocati all'attraversamento dello zero di $\Psi$ anziché alle profondità di pendenza nulla; il picco della curva letto come movimento massimo anziché come inversione del flusso; l'intera curva narrata come un'unica storia continua di velocità dalla superficie al fondale.

\emph{Indicatori negativi} (\emph{non} codificare): inversione di direzione correttamente legata a dove la curva cambia segno di pendenza; trasporto accumulato esplicitamente distinto dalla velocità locale; incertezza produttiva su cosa rappresenti l'altezza della curva --- codificata invece sotto \textsc{uncertainty-as-quality}.

\emph{Nota.} Questo codice è indipendente dal target strutturale della Scena~I. Un partecipante può nominare correttamente le \emph{direzioni} opposte --- ed essere così codificato \emph{grasped} --- pur leggendo ancora l'altezza della curva come \emph{velocità}, e quindi portare anche S1-MISC-INT. I due registrano metà diverse della stessa lettura e non devono essere accorpati: registrarli entrambi è ciò che rivela un partecipante che possiede la direzione del flusso ma non la natura cumulata dello strumento. Incrociata con la condizione, la frequenza del codice è l'evidenza su se la forma esperienziale riduca l'errore; incrociata con la fiducia, su se sia tenuto con falsa certezza.

% -------------------------------------------------------------------------
\subsection{Procedura di codifica}
\label{subsec:app-procedure}

La griglia sopra è stata applicata in due passaggi. Nel primo, a un large language model sono stati
dati la griglia, le sue definizioni di categoria e una risposta alla volta, e ha restituito
un'etichetta proposta. Nel secondo, il ricercatore ha letto ogni risposta rispetto alla griglia e ha
deciso l'etichetta proposta, correggendola ovunque le due letture differissero. Ogni etichetta nel
dataset riportato è quindi del ricercatore; il passaggio automatico è servito a ordinare il lavoro,
non a deciderlo.

Due conseguenze seguono e sono registrate anziché mitigate. Non c'è un secondo codificatore
indipendente, perciò non è riportato alcun coefficiente di accordo tra codificatori --- la misura non
si applica a un disegno a codificatore singolo. E un revisore che vede un'etichetta proposta prima di
formare la propria può ancorarvisi; i criteri fissi della griglia limitano quanto ciò possa andare
avanti, ma non lo rimuovono. Entrambi sono elencati tra i limiti dello studio nella
\S\ref{sec:eval-limitations}.

% -------------------------------------------------------------------------
\subsection{Codici emergenti}
\label{subsec:app-emergent}

Accanto alla griglia definita a priori, la codifica ha registrato pattern che non ne facevano parte.
Questi sono post-hoc: sono emersi durante la lettura, non erano previsti, e sono riportati
separatamente dai target strutturali lungo tutto il Capitolo~\ref{ch:results}. Tre sono ricorsi
abbastanza spesso da essere riportati per condizione; i restanti sono occorsi una o due volte
ciascuno e sono elencati per completezza senza interpretazione, poiché una o due occorrenze non
possono essere distinte dall'idiosincrasia a questa dimensione campionaria.

\begin{description}
\item[\textnormal{\textsc{Interaction-aided}}] La risposta attribuisce la propria comprensione
all'agire sul display --- mettere in pausa, riavvolgere, isolare un elemento. Dodici occorrenze,
tutte nella condizione sperimentale.

\item[\textnormal{\textsc{Static-limit-named}}] La risposta nomina l'assenza di un ordinamento
temporale come la ragione per cui non può dire se il sistema ripassi per stati precedenti. Quattro
occorrenze, tutte nella condizione di controllo.

\item[\textnormal{\textsc{Anomaly-confusion}}] La risposta registra la codifica usuale/insolito ma
non sa dire cosa significhi. Sei occorrenze, una nella condizione di controllo e cinque nella
sperimentale.
\end{description}

I restanti codici, ciascuno occorso una o due volte, sono elencati sotto per completezza. I loro
conteggi si leggono dalla Figura~\ref{fig:emergent-full}.

\begin{description}
\item[\textnormal{\textsc{Diagnosis-by-participant}}] Due occorrenze, entrambe nella condizione di controllo.
\item[\textnormal{\textsc{Mechanism-invented}}] Due occorrenze, entrambe nella condizione sperimentale.
\item[\textnormal{\textsc{Cognitive-load}}] Una occorrenza, nella condizione sperimentale.
\item[\textnormal{\textsc{Color-misread}}] Una occorrenza, nella condizione di controllo.
\item[\textnormal{\textsc{ML-role-named}}] Una occorrenza, nella condizione sperimentale.
\item[\textnormal{\textsc{Prior-knowledge}}] Una occorrenza, nella condizione di controllo.
\item[\textnormal{\textsc{Scene-confusion}}] Una occorrenza, nella condizione sperimentale.
\item[\textnormal{\textsc{Seasonality-misread}}] Una occorrenza, nella condizione di controllo.
\item[\textnormal{\textsc{Sign-confusion}}] Una occorrenza, nella condizione sperimentale.
\end{description}
